
\capitulo{3}{Conceptos teóricos}


Este proyecto se centra principalmente en individuos que parecen un tipo específico de parálisis cerebral, la atetoide. Es necesario comprender los aspectos teóricos detrás de estas afecciones y sus características para poder desarrollar una solución tecnológica óptima y novedosa. 

\section{Conceptos teóricos básicos}
\subsection{¿Qué es la parálisis cerebral?}

La parálisis cerebral es un trastorno neuromotor crónico originado por una lesión o defecto durante el desarrollo del cerebro inmaduro causado durante un periodo comprendido entre los primeros días de gestación y los 5 años de vida \cite{MayoClinic2024}.

La denominación de “parálisis” nos da una gran información sobre las características de la enfermedad, y es que esta genera una gran debilidad y dificultad a la hora de utilizar los músculos, derivando en alteraciones motrices, del tono muscular y de la postura. No solamente afecta a la movilidad, sino que también causa problemas en otros aspectos como la cognición, derivando en déficit intelectual o dificultades en la comunicación \cite{GarciaRon2022, MedlinePlus2024}.

Entre las principales causas que la originan se encuentra la falta de flujo sanguíneo hacia el cerebro en el feto, alteraciones genéticas, infecciones por virus o bacterias, convulsiones o dificultades en el parto \cite{CentrodeNeurologaAvanzada2025}. 

\subsection{Tipos de parálisis cerebral}

Se pueden clasificar en función de la zona del cuerpo afectada:
\begin{itemize}
    \item Hemipléjica: afecta solo a una mitad del cuerpo, o la derecha o la izquierda.
    \item Paraplejia: afecta principalmente a todos los miembros inferiores.
    \item Tetraplejia: afecta a todas las extremidades, tanto a los miembros superiores como inferiores.
    \item Displejia: se ven afectadas las dos piernas, pero las extremidades superiores no se ven nada o casi nada afectadas.
    \item Monoplejia: cuando solo se ve afectado un miembro. 
\end{itemize}

En función de las características que presentan, podemos distinguir las siguientes:
\begin{itemize}
    \item Parálisis cerebral espástica: es la más común dándose entre un 60-70\% de los pacientes que padecen parálisis cerebral. Provoca rigidez y dificultad para controlar los músculos \cite{ASPACE, ASPACEGalicia2024}. 
    \item Parálisis cerebral atetoide: que provoca movimientos lentos, involuntarios y descoordinados. Deriva en problemas para controlar el movimiento de extremidades y en dificultades para estar sentado o caminando \cite{ASPACE, MedlinePlus2024}.
    \item Parálisis cerebral atáxica: afecta la coordinación y el equilibrio derivando en posturas y movimientos inestables \cite{eFISIO2023}. A pesar de las dificultades, existen personas que pueden llegar a caminar, aunque con muchas dificultades.
    \item Parálisis cerebral mixta: afecta a diferentes áreas cerebrales provocando una combinación de las dificultades mencionadas anteriormente en los otros tipos existentes.
\end{itemize}

Otra forma de clasificación es en función de si es leve, cuando se presentan alteraciones, pero no se ven limitadas las actividades diarias, moderada, cuando sí que existen dificultades en las actividades diarias y necesitan asistencia y apoyo, severa, cuando requiere ayuda para todas las actividades \cite{ASPACE}. También se pueden diferenciar en función del momento en el que aparece la lesión en prenatal, perinatal o postnatal \cite{ValldHebron2026}.


\subsection{Parálisis cerebral ateoitde}

Para este proyecto interesa la parálisis cerebral atetoide debido a sus característicos movimientos involuntarios, por lo que se va a profundizar en esta para comprender mejor con qué tipo de pacientes se trata.

También conocida como discinética, es el segundo tipo de parálisis más común presentándose entre el 10-15\% de personas con parálisis y requiere de terapia y asistencia de por vida \cite{PercyMartnez2025}. Se caracteriza por ser trastorno permanente no progresivo que se origina por una lesión cerebral que tiene lugar durante el desarrollo fetal o a lo largo de la primera infancia \cite{Li2022}. Su principal característica es que dificulta el control y coordinación muscular provocando movimientos involuntarios, que suelen afectar a los brazos las piernas o el rostro afectando a el desarrollo y la calidad de vida del paciente \cite{eFISIO2022}.

Su aparición se suele relacionar con los ganglios basales y lesiones talámicas que además pueden estar acompañadas de periodos hipóxicos breves. Una de las causas es kernicterus derivado de hiperbilirrubinemia neonatal provocando que la bilirrubina se acumule en los ganglios basales. Hoy en día la hiperbilirrubinemia se previene y trata con mayor eficacia por lo que ha disminuido la presencia de esta enfermedad, pero, aun así, sigue presente por vías como hemorragias intracraneales, ictus o infecciones de los ganglios basales o tálamo.

El tratamiento está enfocado en el control de los síntomas y la mejora de la calidad de vida del paciente. Para esto se emplean medicamentos para la reducción del dolor en caso de que haya y medidas para el tratamiento de los síntomas generalmente mediante la asistencia de un equipo integrado por profesionales clínicos, terapeutas y especialistas en salud conductual que tratan de maximizar la calidad de vida.

También existen medicamentos como el GABA-B, la estimulación cerebral profunda o la toxina botulínica que se pueden emplear pero que hasta el momento no han demostrado tener una gran eficacia \cite{Li2022}.

\section{Estado del arte y trabajos relacionados.}

Al llevar a cabo un proyecto en el que se aporta una solución tecnológica, es muy importante tener en cuenta qué se ha desarrollado hasta el momento para no repetirlo. El objetivo es aportar novedades que o bien creen soluciones completamente nuevas o mejoren las ya existentes.

\subsection{Wearable Devices for Assessment of Tremor}

Este artículo realiza una revisión sobre el creciente uso de los dispositivos wereables a la hora de evaluar el temblor en determinadas enfermedades. Se analizan los sensores inerciales, giroscopios y acelerómetros como medio de evaluar el movimiento de manera que puedan ser útiles a la hora de valorar la progresión de un trastorno como el Parkinson. 

Concluye que existen desafíos importantes a la hora de distinguir el temblor de otros movimientos involuntarios pero que este tipo de dispositivos tienen un gran potencial. \cite{Vescio2021}

\subsection{Wearable inertial sensors for human movement analysis: a five-year update}

Evalúa el empleo de sensores inerciales durante los últimos años como métodos de seguimiento cotidiano y prolongado. Presentan un gran potencial a la hora de analizar enfermedades neurológicas que generen alteraciones motoras por su capacidad para detectar patrones de movimiento. 

Su mayor reto está en ser capaces de elaborar algoritmos que interpreten de manera adecuada las señales que captan para poder dar una interpretación clínica fiable. \cite{Picerno2021}
